% Vista científica exacta materializada desde una rama condicional.
% Fuente: source/demostraciones_primarias/14_05b_extension_superviviente_caracter.tex; rama: HMTIncludeRouteLedgerBlock.
% No modifica el contenido matemático de la rama de origen.
\chapter{Extension survivante, registre dodécaphasique et signature entière}
\label{chap:extension-ruta-caracter}

La trajectoire observée ne constitue pas l'objet premier de cette
construction. L'objet premier est une extension compatible à toutes les
profondeurs du transport HMT ; un chemin fini est l'une de ses lectures. Cette
distinction permet de formuler sans ambiguïté la chaîne interne qui conduit d'une
histoire enrichie au caractère d'action. Elle détermine aussi quelles
informations peuvent être oubliées sans modifier la sortie et lesquelles doivent rester dans le
registre.

Le but de cette section est de réunir en un seul argument quatre éléments qui
avaient été étudiés séparément : l'extension survivante, la mémoire
intrinsèque, le registre orienté et la signature entière. Le résultat terminal de
cette section est le caractère formel
\[
 \Gamma_{108}^{\rm enr}
 \xrightarrow{\ L_{\rm tick}\ }
 \ZZ^5
 \xrightarrow{\ \chi_{\rm form}\ }
 \operatorname{Fun}(\mathcal P,\RR_{>0}),
\]
antérieur tant à la spécialisation des coefficients qu'à toute unité
physique. La première application extrait la signature du registre lu ; la seconde la
convertit en une fonction positive sur l'espace des paramètres
\(\mathcal P\). La spécialisation numérique ne s'effectue qu'après
la construction des coefficients globaux dans la section correspondante. La
réalisation métrologique relève d'une carte dimensionnelle encore ultérieure.
L'application \(L_{\rm tick}\) est le lecteur entier d'itération et de retour du
registre chronologique enrichi.

\section{Composition opératoire de la prolongation et cochaîne torsionnelle dans la génération du registre de chemin}

La loi des masses commence avant la formule exponentielle. APP apporte un unique
support de chemins et ses deux feuillets, somme et produit. Le TRIT oriente chaque
itération élémentaire.
Si
\[
 r_t=1+((t-1)\bmod9),\qquad m_t:=M_{\rm ph}(r_t),
\]
le sélecteur de phase \(M_{\rm ph}\) active
\[
 \operatorname{op}_t=
 \begin{cases}
  \Sigma,&m_t=+1,\\
  \Pi,&m_t=-1,\\
  \varnothing\quad(\text{\rm rétention ou compteur neutre}),&m_t=0.
 \end{cases}
\]
La lecture temporelle peut interpréter le troisième cas comme seuil ou présent,
mais ne le transforme pas en un troisième feuillet arithmétique. Ce n'est pas APP qui alterne
ses feuillets. TPK transporte
conjointement position, orientation, feuillet, phase, bloc, retenue, préfixe,
cylindre, généalogie et mémoire. Le registre chronologique enrichi note
chaque transition ---position,
feuillet, lecture, retenue, phase, bloc et événement---. La signature locale de frontière
est ensuite dérivée au moyen de \(\operatorname{Sig}_q(h,c,\lambda)\), et la signature
globale de masse est extraite du registre complet ; aucune n'est reconstruite
rétrospectivement à partir d'une masse.

L'entonnoir fini qui alimente cette histoire conserve quatre types distincts :
\[
 9^2\!\cdot9^2\!\cdot4\!\cdot4=104\,976
 \longrightarrow 468\ \text{émissions }U_6
 \longrightarrow 243\ \text{mots }w_6
 \hookrightarrow \mathbb F_3^6,
 \qquad |\mathbb F_3^6|=729.
\]
Le premier chiffre compte les conditions initiales du moteur à deux curseurs ; 468
compte les émissions décimales et 243 les mots ternaires visibles. Les 324 germes
de la projection chronologique \(\Pi_{108}^{\rm cal}\) à un seul curseur
ne remplacent pas non plus les 104\,976
conditions de l'état complet. Chaque flèche change donc de domaine et
d'opération. L'égalité numérique entre l'espace ambiant de \(729\) mots et les
\(729\) suffixes d'une étape de filtration n'est reliée que par le codage
ternaire déclaré ; elle n'identifie pas les deux objets.

Le relèvement conserve également son ordre causal :
\[
 w_6\xrightarrow{\,L_0\,}w_{12}
 \xrightarrow{\,L_0\,}w_{18}
 \xrightarrow{\,L_1\,}w_{24}
 \xrightarrow{\,L_1\,}w_{30}
 \longrightarrow R_{36}\longrightarrow\mathcal G_9
 \longrightarrow
 \text{section survivante}.
\]
La frontière \(R_{36}\) ouvre la prolongation. À chaque étape de filtration,
\(\mathcal G_9\) épuise
les 729 sous-cylindres et conserve l'unique compatible avec tous les horizons ;
l'épisode spéculaire profond \(3\to1\) est un autre élagage et ne doit pas être confondu avec
cette partition ordinaire \(729\to1\). Le retour nonadique conserve la phase,
mais change le bloc, le cylindre et la mémoire. Ni chiffres cibles ni masses n'interviennent
dans cette sélection. À la frontière dimensionnelle, une particule est une section
persistante enrichie de l'extension ; le chemin est son histoire observable.
Son nom physique se lit ensuite et ne décide pas quelle branche survit.

Avant d'ouvrir le registre chronologique dimensionnel, la même généalogie conserve une
fenêtre publiée de vingt sextuors par canal et admet la première lecture
archimédienne de douze triades décimales. Ces deux témoins appartiennent à la
prolongation de l'état : ils ne sont pas les douze événements du registre de masse et ne les
remplacent pas. Les événements ne se construisent qu'ensuite, en lisant le chemin sur le
système dodécaphasique orienté.

La généalogie dodécaphasique convertit ensuite les douze événements orientés en un
registre et seulement alors en la signature
\[
 \gamma_x\longmapsto\mathcal L(\gamma_x),\qquad
 \sigma_{\rm reg}:=L_{\rm tick}(\gamma_x)
 =(n_A,s_C,k_{\Dfour},b_{120},b_{\zeta}),
\]
\[
 \sigma_{\rm reg}\longmapsto
 \chi_{\rm form}(\sigma_{\rm reg};\,\cdot).
\]
Les coordonnées publiques \(b_{120}:=\nu_{120}^{\rm ev}\) et
\(b_{\zeta}:=\nu_{270}\) sont des sorties du registre chronologique, non des corrections
métrologiques introduites à la fin. Les coefficients ultérieurs du
semi-groupe de tours conservent une autre provenance :
\[
 N_{120}=b_{120}+\eta_{120},\qquad
 b_{\zeta}\zeta_{270}\not\equiv\eta_{270}s_{270},
\]
avec \(\zeta_{270}=135\Dfour^2\) et
\(s_{270}=q_-^{270}-q_+^{270}\). Le raffinement des tours s'ajoute après
le caractère central sans effacer la provenance d'aucune coordonnée.

Trois sorties latérales conservent aussi leur type. Le secteur nul quitte la
branche positive au moyen de son lecteur et de sa carte nuls ; il n'est pas la limite
\(m\to0\) d'un caractère positif. Une trajectoire virtuelle est une extension
finie avec obstruction terminale et conserve le registre chronologique et la retenue sans acquérir pour
cela une masse imaginaire. Enfin, l'ombre de Witt accompagne chaque étape de filtration et
préserve la généalogie exceptionnelle, mais ne sélectionne ni le chemin ni les
coefficients de masse.

\begin{figure}[p]
\centering
\begin{tikzpicture}[
  font=\small,
  node distance=3.5mm,
  box/.style={draw,rounded corners=1.5mm,align=center,text width=12.8cm,
    inner xsep=6pt,inner ysep=4pt},
  hmt/.style={box,draw=hmtblue,fill=hmtlightblue},
  md/.style={box,draw=hmtviolet,fill=hmtlightviolet},
  act/.style={box,draw=hmtgreen,fill=hmtlightgreen},
  result/.style={box,draw=hmtgold,fill=hmtlightgold},
  arr/.style={-{Stealth[length=2.2mm]},line width=.7pt}]

  \node[hmt] (app) {\textbf{APP} : tore arithmétique nonadique, chemins et feuillets additif et multiplicatif. La spécification complète contient \(104\,976\) graines à double curseur.};
  \node[hmt,below=of app] (trit) {\textbf{TRIT} : orientation locale de chaque transition ; le relèvement conserve le résidu, le quotient et la retenue.};
  \node[hmt,below=of trit] (tpk) {\textbf{TPK} : transport de phase avec mémoire ; \(104\,976\to468\,U_6\to243\,w_6\subset\mathbb F_3^6\).};
  \node[hmt,below=of tpk] (surv) {Filtration coinductive et monodromie \(\mathcal G_9\) : trajectoire survivante à généalogie conservée.};
  \node[md,below=of surv] (route) {La trajectoire de particule produit, sans consulter la masse, un enregistrement dodécaphasique et la signature \(\sigma=(n_A,s_C,k_{\Delta_4},b_{120},b_\zeta)\in\mathbb Z^5\).};
  \node[md,below=of route] (chi) {Le caractère central \(\chi_0(\sigma)\) et le raffinement harmonique \(\eta\in\mathbb Z^{\langle90,120\rangle}\) construisent \(\chi_{\rm full}(\sigma,\eta)\in\mathbb R_{>0}\).};
  \node[act,below=of chi] (action) {Branche dimensionnelle indépendante : \(\Sigma_H=\varphi\alpha_{\rm HMT}^{16}\to\hbar_{\rm int}\), \(\omega_0=2\pi/(108t_0)\), \(m_0=\hbar_{\rm int}\omega_0c_{\rm int}^{-2}\).};
  \node[result,below=of action] (mass) {Confluence typée : \(m_X^{\rm int}=m_0\,\chi_{\rm full}(\sigma_X,\eta_X)\). L’échelle commune fixe la normalisation absolue et s’annule dans \(m_X/m_Y\).};
  \node[result,below=of mass] (compare) {La carte d’unités exprime la grandeur. CODATA et PDG ne sont consultés que dans la dernière colonne de comparaison.};

  \foreach \a/\b in {app/trit,trit/tpk,tpk/surv,surv/route,route/chi,
    chi/action,action/mass,mass/compare}{\draw[arr] (\a)--(\b);}
\end{tikzpicture}

\caption{Dérivation d'une grandeur de masse à partir d'une trajectoire HMT.
Le caractère apporte un rapport adimensionnel ; l'échelle d'action et l'horloge
apportent la section dimensionnelle. La référence externe n'intervient qu'après
l'application de la carte d'unités.}
\label{fig:derivacion-masa-rev9}
\end{figure}

\section{Projection d'une extension survivante sur son chemin observable}

Soit \((\mathcal S_m,\rho_{m+1,m})_{m\geq0}\) le système d'états HMT
survivants à la profondeur \(m\), avec des applications de troncature
\(\rho_{m+1,m}:\mathcal S_{m+1}\to\mathcal S_m\). Nous définissons
\[
 \Omega_{\HMT}^{\rm surv}
 :=\varprojlim_m\mathcal S_m
 =\left\{(x_m)_m:\rho_{m+1,m}(x_{m+1})=x_m\right\}.
\]
Un élément \(x\in\Omega_{\HMT}^{\rm surv}\) conserve simultanément la
compatibilité de toutes ses troncatures. Une carte de lecture de profondeur
108 produit
\[
 \operatorname{sh}_{108}(x)=\gamma_x\in\Gamma_{108}^{\rm enr}.
\]
La notation \(\operatorname{sh}\) rappelle que \(\gamma_x\) est une ombre
observable : elle enregistre une présentation séquentielle de l'état, mais l'état ne
se définit pas en choisissant rétrospectivement ce chemin.

La présentation enrichie conserve, outre le mot visible, la phase,
le feuillet, l'orientation, le bord, la mémoire antifeuillet et les prédécesseurs nécessaires pour
évaluer les cocycles. Deux mots ayant les mêmes cinq entiers finaux peuvent
être des présentations distinctes parce que ces entiers sont calculés après
l'application des signes et des corrélations.

\begin{definition}[Mémoire intrinsèque et registre lu]
Une extension dodécaphasique enrichie est une classe compatible de
troncatures. Nous notons \(\mathscr M(x)\) la mémoire intrinsèque que ces
troncatures conservent. La carte
\(\gamma_x=\operatorname{sh}_{108}(x)\) étant fixée, son registre lu est
\[
 \mathcal L(\gamma_x)
 =\bigl(n_A,e_0,\ldots,e_{11},T_\zeta,C_\zeta\bigr),
\]
où \(n_A\) est le dénombrement du support d'action déclaré, les \(e_m\) sont
les douze événements orientés et \(T_\zeta,C_\zeta\) sont la mémoire torsionnelle
totale et sa restriction à la couronne. La compatibilité exige que ces données
soient préservées par toute troncature contenant les pas employés dans leur
définition.
\end{definition}

L'extension contient déjà la mémoire que le chemin présente ; le chemin ne la crée pas.
Relativement à la carte déclarée, la composition
\[
x\longmapsto\gamma_x
\longmapsto\mathcal L(\gamma_x)
\longmapsto L_{\rm tick}(\gamma_x)
\]
est une application. L'invariance par rapport à une autre carte de présentation
exigerait en outre de démontrer la compatibilité des applications
de transition correspondantes. En ce sens précis, le registre ne s'ajoute pas à une
particule déjà construite : c'est une lecture de la mémoire de l'extension.

\section{Décomposition entière de \(12\) événements dans les secteurs \(1+5+6\)}

Soit \(d_1,\ldots,d_{108}\in\{1,\ldots,9\}\) le mot d'une présentation
enrichie. Dans la convention de parité utilisée par l'extracteur, le support
d'action est
\[
 \mathcal D_{\rm act}=\{1,3,5,7\},
 \qquad
 n_A=\#\{t:d_t\in\mathcal D_{\rm act}\}.
\]
Le choix de ce support est une partie déclarée de la carte ; il ne doit pas
être confondu avec la classification TRIT des pas. Pour les événements pairs, on
utilise
\[
 \mathcal D_{\rm evt}=\{2,4,6,8\},
 \qquad
 \Sigma_{12}=(+,+,+,-,-,-,+,+,+,-,-,-).
\]
L'événement orienté du bloc \(m\in\{0,\ldots,11\}\) est
\[
 e_m=\Sigma_{12}(m)
 \#\{t\in\{9m+1,\ldots,9m+9\}:d_t\in\mathcal D_{\rm evt}\}.
\]

Apparions les deux demi-couronnes au moyen de
\[
 p_i=e_i+e_{i+6},
 \qquad
 a_i=e_i-e_{i+6},
 \qquad 0\leq i<6.
\]
On définit
\[
 s_C=\sum_{i=0}^{5}p_i,
 \qquad
 M_i=p_i-p_5\quad(0\leq i<5),
\]
et nous écrivons \(M_5=(M_0,\ldots,M_4)\),
\(A_6=(a_0,\ldots,a_5)\).

Le \cref{thm:memoria-1-5-6}, dans la résidence primaire du système
dodécaphasique de TPK, démontre que
\[
 (e_0,\ldots,e_{11})\longmapsto(s_C,M_5,A_6)
\]
est injectif sur son image et fournit l'inverse
\[
 p_5=\frac{s_C-\sum_{i=0}^{4}M_i}{6},\qquad
 p_i=p_5+M_i,
\]
\[
 e_i=\frac{p_i+a_i}{2},\qquad
 e_{i+6}=\frac{p_i-a_i}{2}.
\]
On ne récupère ici que ces coordonnées pour construire le caractère ; la
preuve et l'identité \(1+5+6=12\) ne sont pas dupliquées.

La donnée \(s_C\) est celle qui entre dans le caractère angulaire. Les onze coordonnées
\((M_5,A_6)\) conservent le placement de ce total. Les éliminer avant de
composer des histoires confond des événements différents susceptibles de produire
des corrélations distinctes.

\section{Cochaîne torsionnelle au niveau du pas}

Le dénombrement des événements ne détermine pas à lui seul la contribution torsionnelle. Pour
chaque occurrence de neuf, nous conservons le chiffre qui la précède et définissons
\[
 j_2(d)=
 \begin{cases}
 0,&d=9,\\
 +1,&d\in\{2,4,6,8\},\\
 -1,&d\in\{1,3,5,7\}.
 \end{cases}
\]
Avec des indices cycliques dans le mot, soit
\[
 \zeta_t=\mathbf 1_{\{d_t=9\}}j_2(d_{t-1}),
\]
\[
 T_\zeta=\sum_{t=1}^{108}\zeta_t,
 \qquad
 C_\zeta=\sum_{t\in\{27,54,81,108\}}\zeta_t.
\]
La coordonnée forte qui accompagne \(\Dfour\) est
\[
 \boxed{k_{\Dfour}=T_\zeta-2C_\zeta.}
\]
La formule sépare l'entraînement accumulé de sa correction orientée dans les
quatre clôtures de couronne. Le recensement exhaustif montre que ni
\(T_\zeta\) ni \(C_\zeta\), séparément, ne séparent toutes les fibres du recensement ;
le couple ordonné le fait dans les deux conventions typées qui sont confrontées.

Les deux coordonnées de tour restantes sont calculées à partir du registre
d'événements. Soit
\[
 \tau_m=\operatorname{sgn}(e_m).
\]
Pour les quatre classes de pas quatre,
\[
 I_a=\{a,a+4,a+8\}\pmod{12},
 \qquad 0\leq a<4,
\]
nous définissons le caractère pondéré par les événements
\[
 \nu_{120}^{\rm ev}
 =\sum_{a=0}^{3}\operatorname{sgn}
 \left(\sum_{m\in I_a}e_m\right),
\]
et l'autocorrélation de pas trois
\[
 \nu_{270}=\sum_{m=0}^{11}\tau_m\tau_{(m+3)\bmod12}.
\]
La variante
\[
 \nu_{120}^{\rm sgn}
 =\sum_{a=0}^{3}\operatorname{sgn}
 \left(\sum_{m\in I_a}\tau_m\right)
\]
est une autre observable. Elles ne sont pas identifiées tacitement : elles diffèrent sur
25 des 81 chemins de l'atlas.

\begin{definition}[Extracteur enrichi de signature]
La convention d'événements étant fixée, on définit
\[
 L_{\rm tick}(\gamma)=
 \bigl(n_A,s_C,k_{\Dfour},\nu_{120}^{\rm ev},\nu_{270}\bigr)
 \in\ZZ^5.
\]
L'indice indique que la coordonnée torsionnelle est calculée avant d'oublier
les prédécesseurs des pas nuls.
\end{definition}

Cette application est déterministe et ne consulte ni masses, ni unités, ni étiquettes
de particules. On n'affirme pas qu'elle soit linéaire par rapport à une concaténation déjà
projetée : les signes et les autocorrélations imposent de composer d'abord les
registres enrichis.
